\documentclass{article}
\usepackage[T2A]{fontenc}
\usepackage[utf8]{inputenc}
\usepackage{amsthm}
\usepackage{amsmath}
\usepackage{amssymb}
\usepackage{amsfonts}
\usepackage{mathrsfs}
\usepackage[12pt]{extsizes}
\usepackage{fancyvrb}
\usepackage{indentfirst}
\usepackage[
  left=2cm, right=2cm, top=2cm, bottom=2cm, headsep=0.2cm, footskip=0.6cm, bindingoffset=0cm
]{geometry}
\usepackage[english,russian]{babel}


\begin{document}
\section*{Вариант 22}
Рассмотрим дискретную временную шкалу $T=\{t_0<t_1<\ldots<t_N\}$, в узлах которой 
заданы значения двух показателей: $\{q(t_k)\}$ (например, объёма спроса, ВВП и пр.), и
$\{p(t_k)\}$ (цены, объёма трудовых или капитальных ресурсов и пр.), $k=0,\ldots,N$. Считаем, что
$\min\limits_{k\in0,\ldots,N} p(t_k)=p(t_0)<\min\limits_{k\in1,\ldots,N} p(t_k)$ или
$\max\limits_{k\in0,\ldots,N} p(t_k)=p(t_0)>\max\limits_{k\in1,\ldots,N} p(t_k)$,
$\min\limits_{k\in0,\ldots,N} q(t_k)=q(t_0)<\min\limits_{k\in1,\ldots,N} q(t_k)$ или
$\max\limits_{k\in0,\ldots,N} q(t_k)=q(t_0)>\max\limits_{k\in1,\ldots,N} q(t_k)$.

Введём в рассмотрение нормированные индексы, соответственно, для $k=0,\ldots,N$, взяв в
качестве базы индексирования начальные значения показателя в точке $t_0$:
\begin{align}
I_q(t_k)=q(t_k)/q(t_0), \\
I_p(t_k)=p(t_k)/p(t_0).
\end{align}

Вычислим точечную эластичность показателя $q$ относительно показателя $p$ для каждого
$k=0,\ldots,N$ по формуле:
\begin{align}
E(t_k)=\frac{I_q(t_k)-1}{I_p(t_k)-1}=\frac{\Delta q/q}{\Delta p/p}.
\end{align}

\end{document}

